% Generated draft: requires mathematical and source review.
\documentclass[11pt]{article}
\usepackage[margin=25mm]{geometry}
\usepackage{fontspec}
\setmainfont{Latin Modern Roman}
\usepackage{amsmath,amssymb,mathtools,unicode-math}
\setmathfont{Latin Modern Math}
\usepackage{enumitem,needspace,xurl,hyperref}
\setlength{\parindent}{0pt}
\setlength{\parskip}{5pt}
\setlength{\emergencystretch}{3em}
\newcommand{\N}{\mathbb{N}}
\newcommand{\Z}{\mathbb{Z}}
\newcommand{\Q}{\mathbb{Q}}
\newcommand{\R}{\mathbb{R}}
\newcommand{\C}{\mathbb{C}}
\newcommand{\F}{\mathbb{F}}
\newcommand{\A}{\mathbb{A}}
\newcommand{\PP}{\mathbb{P}}
\newcommand{\E}{\mathbb{E}}
\newcommand{\eps}{\varepsilon}
\newcommand{\dd}{\,\mathrm{d}}
\begin{document}
\section*{UnsolvedMath ID: 30000526}
\textbf{Source problem code:} catalogue entry
\subsection{Definitions and mathematical statement}
On the domain intended by the report, consider
\[
 \operatorname{div}(|Du|^{p-2}Du)
     =\operatorname{div}(|F|^{p-2}F).
\]
In the usual scalar elliptic interpretation, $\Omega\subseteq\R^n$ is open, $p>1$, $u$ is scalar valued, $F:\Omega\to\R^n$, and $Du$ is the distributional gradient. These domain, dimension, parameter, and solution-class conventions are not explicitly supplied.

When the two fluxes are locally integrable, the displayed equation means
\[
 \int_\Omega |Du|^{p-2}Du\cdot D\varphi\,\dd x
  =\int_\Omega |F|^{p-2}F\cdot D\varphi\,\dd x
 \qquad(\varphi\in C_c^\infty(\Omega)).
\]
The nonlinear flux at zero is interpreted as zero. For a finite exponent $q>0$, membership in $L^q_{\mathrm{loc}}$ means finite integral of the $q$-th power on every compactly contained subdomain, including the quasi-norm interpretation if $q<1$.

The requested implication is
\[
 F\in L^q_{\mathrm{loc}}(\Omega)
       \quad\Longrightarrow\quad
 Du\in L^q_{\mathrm{loc}}(\Omega)
 \qquad\text{throughout the source's full range }q\ge p-1.
\]
The lower endpoint $q=p-1$ is explicitly included in the source and is not silently replaced by $q>p-1$, weak-$L^q$, or a smaller open range. Whether the range includes $q=\infty$ is not stated; no infinite-exponent assertion is added. The supplied statement asks for strong local membership, not an explicit norm estimate or boundary regularity result.

The admissible solution notion is crucial. Energy solutions in $W^{1,p}_{\mathrm{loc}}$ and very weak or other distributional solutions have different built-in integrability, especially below $p$. The source gives no a priori gradient integrability, admissibility or approximation condition, scalar-versus-system convention, or restrictions on $p,n$. These must be recovered before asserting the implication for every distributional solution. The cited parabolic extension is contextual and does not automatically settle this elliptic equation.
\subsection{Short English statement}
The question asks whether the gradient inherits the same strong local integrability as the data throughout the stated exponent range, including $q=p-1$. The solution class and endpoint conventions are missing, and no weak-space or parabolic substitute is imposed.
\subsection{Sources}
\begin{itemize}
\item Local UnsolvedMath registry, permanent ID 30000526, code catalogue entry, ``Full-Range Calder\'on--Zygmund Estimates for the $p$-Laplacian.'' Proposed year 2006; proposer unspecified.
\item \emph{Calculus of Variations} (2007), Oberwolfach Reports, pp.\ 1879--1940, DOI 10.4171/owr/2006/31; \url{https://doi.org/10.4171/owr/2006/31}. The displayed report year 2007 differs from the DOI and proposed year 2006.
\item Supplied reference ``parabolic Calderon--Zygmund extension'': \url{https://arxiv.org/abs/2109.02595}.
\item Statement status corrected\_verified; the supplied verification says official-report content was checked and notation, context, and question form normalized without changing the source claim.
\item Literature triage checked 2026-08-21: open. It says the report leaves the full $q\ge p-1$ range open and later results extend portions without establishing the extracted general statement. Precise solution classes, endpoint claims, and current status require independent verification. Created 2026-08-25; updated 2026-09-16.
\item Registry curation by Ulam AI / UnsolvedMath Contributors, CC BY 4.0; underlying report and literature provenance remain as supplied. Unreviewed statement-only draft.
\end{itemize}
\end{document}
